\section{Original local line coefficients compute genuine model section norms}

\begin{lemma}[The original global section has its genuine whole local scalar morphism]
\label{lem:fc-original-local-section-scalar}
\leanok
\lean{EnvelopeProperty.SchemeLineSectionEvaluation.pullbackSectionHom_restrict, EnvelopeProperty.SchemeLineSectionEvaluation.sectionHom_localScalar}
\uses{lem:fc-original-line-section-frame-naturality}
Restrict an actual original global module section to a genuine affine open
and apply any actual original line frame.  The pullback of its original
structure-line section morphism, followed by the actual pullback frame, is
exactly multiplication by the actual original local frame coefficient.
No scalar-compatibility certificate is supplied for this original section.
\end{lemma}
\begin{proof}
The naturality square for the native restriction--pullback isomorphism
identifies pullback of the original section map with its native restriction.
The genuine structure-unit comparison supplies the source identity.
The original global-section/unit-Hom equivalence commutes with restriction
and with composition by the given frame.  Its coefficient is the original
local trivialization top section.  The actual affine section transport and
its inverse identify this with the original ambient local coefficient.
\end{proof}

\begin{lemma}[An actual original local scalar evaluates at the same original point]
\label{lem:fc-original-local-section-residue-value}
\leanok
\lean{EnvelopeProperty.Berkovich.RelativeSchemeValuedPoint.localGenericFactor_openTopSection}
At the genuine proper center of a prescribed-base valued point, the actual
local generic factor pulls each original local section to exactly its
original residue-field evaluation.  This is an equality in the SAME actual
original residue field.
\end{lemma}
\begin{proof}
Transport the original section to the top open of its genuine open subscheme.
This is exactly the actual open-inclusion section map.  Native section-map
composition and the whole local-generic factor triangle identify the
result with the original residue-field pullback.  The genuine native
residue-field evaluation theorem identifies that pullback with the original
scheme evaluation.  Quantify the open-containment proof before using the
whole scheme triangle, so the identity retains the original local section.
\end{proof}

\begin{theorem}[The genuine original model norm is the original local coefficient value]
\label{thm:fc-original-local-section-model-norm}
\leanok
\lean{EnvelopeProperty.Berkovich.RelativeSchemeValuedPoint.modelLineNorm_fibreSection_eq_localScalar, EnvelopeProperty.Berkovich.RelativeSchemeValuedPoint.modelLineNorm_fibreSection_eq_localEval}
\uses{lem:fc-original-line-section-fibre,lem:fc-original-local-section-residue-value}
For a genuine invertible original module on the actual proper model, an
actual original local frame at the genuine center computes the model norm
of the SAME original global-section fibre as the original valued-point
evaluation of its SAME original local coefficient.  An internal whole
local scalar-morphism identity suffices, and the preceding section theorem
derives that identity for actual original global module sections.
\end{theorem}
\begin{proof}
The genuine local-frame model norm theorem computes the actual norm by the
original residue valuation of the actual frame coefficient.  Apply the
whole pullback scalar identity to one at the abstract scheme level, and
then apply the residue valuation before specialization to the actual
residue-field spectrum.  This gives the first formula with the actual
pullback scalar.  The preceding original residue-field evaluation identity
then gives the original local valued-point evaluation formula.  No norm
comparison is an input to the resulting theorem: the frame norm formula
comes from the proved proper valuative-frame construction.
\end{proof}
